\section{Conclusion}
\label{sec:conclusion}

\emph{RQ1.} AI Act obligations can be translated into evidence records through four steps, from the legal text to a pipeline requirement, to the fields of a record and to the mechanism that makes it verifiable, provided each derivation states what the record does not establish. Applied to seven provisions, the method yields a matrix whose last column bounds every claim, and an evidence model whose decisions re-execute from the record: all 1,788 regenerated records and all 16 records of the credit-scoring case reproduced their decision and rule.

\emph{RQ2.} Signed, hash-chained records let a third party detect any alteration of a record after issuance. An external anchor adds truncation and rewriting by a key holder, up to the last anchor. Anchoring every record on a permissioned ledger closes that window and adds the identity of the submitting organisation. No configuration detects an event withheld or forged before it reaches the engine, except a forged approval, which the engine rejects. Redaction for Art.~78 kept every integrity anchor and disclosed none of the 60 items checked, once a keyed commitment replaced the policy digest, from which the policy could otherwise be recovered by search. Pseudonymisation remains reversible by whoever holds its key.

\emph{RQ3.} The cryptographic cost is small. A verifiable log adds 0.068\,ms per event to a database row, of which the signature takes 0.017\,ms, and the cost does not grow with the ledger. The cost of a ledger lies in ordering: with serial submission each transaction waits for the orderer's batch timeout, 2,028\,ms at 2\,s and 278\,ms at 250\,ms. Batching raises durable-evidence throughput from 0.49 to 245 records per second at the default timeout and to 1,730 at 250\,ms, at the price of the time a record waits for its batch.

AI Act obligations for high-risk systems can thus be turned into third-party-verifiable evidence produced as a by-product of the ML pipeline. The cryptographic cost is small, and the dominant cost of ledger anchoring is distributed ordering. The evidence establishes the integrity and provenance of recorded decisions. Their truth and their legal sufficiency lie outside it.

Future work runs along five lines:
\begin{enumerate*}[label=(\roman*)]
\item a naturalistic evaluation with auditors, notified bodies and authorities;
\item a production deployment with several orderers, TLS, hardware security modules, qualified seals and time stamps, and the qualification of the ledger under eIDAS;
\item authentication at the source, through signed producers and an attested engine;
\item agentic AI systems, whose decisions are harder to delimit as events; and
\item post-quantum migration, with policies for re-signing and time-stamping existing evidence.
\end{enumerate*}
